\section{Three ways to house a colony}
\label{sec:habitats}

The model represents three habitat strategies. They differ in where the
volume comes from, and it is that difference --- rather than any difference in
comfort or technology --- that the model is built to price.

\subsection{Surface construction}

Ground is zoned and the settlement raises buildings on it as demand and land
value allow, to a density ceiling set by the current era. A high-density
habitation tile at its final stage houses 62 residents. Surface construction
is available essentially everywhere: across all site classes tested, around
$16{,}100$ of the $16{,}384$ tiles are buildable. It is constrained not by
where it may go but by land value, by the era, and by the fact that every
zoned tile must be within one tile of all three networks.

\subsection{Bored arcologies}

Four structures excavate downward into permanently shadowed ice. They may only
be sited on an ice deposit, which --- because of the coupling described in
Section~\ref{sec:model} --- confines them to shadowed crater floors, and they
require a level pad of $3\times3$ or $5\times5$ tiles.

Each opens as a single gallery and sinks one level at a time, but only while
all three networks reach it and the colony is neither short of generating
capacity nor of pressurisation. This is the same condition under which zoned
ground develops, applied to a structure. A bore that loses its utilities stops
where it is and keeps what it has already opened.

Their yield depends on the ice within reach, summed over a radius of three
tiles and \emph{shared} with any other bore whose radius overlaps. Two bores
placed close together therefore make each other worse, which makes the
distribution of the deposit --- not merely the presence of one tile of it ---
the object of the siting decision.

The four differ in what they return. One is a pure housing structure; one
melts its ice and returns a large part of its own pressurisation; one refines
and exports, employing more people than it houses and venting dust that
degrades the ground downwind; and one, gated behind the final era, is the
largest structure in the model.

\subsection{Lava tube arcologies}

The third strategy takes a void that already exists. Where the terrain
generator carves a sinuous rille, the model records the intact tube continuing
beneath it: a course of tiles, a width, and a roof depth. A tube arcology is
sunk onto that course and then extends \emph{along} it, one tile at a time,
under the same utility condition as a bore.

Three properties distinguish it, and together they are the reason it is
modelled as its own family rather than as a fifth bore:

\begin{enumerate}
\item \textbf{The volume is free.} Nothing is excavated, only sealed, lit and
  pressurised. It is accordingly the cheapest arcology per resident in the
  model.
\item \textbf{The ceiling belongs to the map.} A bore descends to a depth fixed
  in its own data; a tube arcology stops when it runs out of tube. Across
  generated polar worlds the usable tube runs from 24 to 81 tiles against a
  structural reach limit of 64, so on some maps the structure binds and on
  others the geology does. The player cannot lengthen a tube.
\item \textbf{It returns no atmosphere.} It has no ice to crack, so every one
  of its residents breathes oxygen produced elsewhere in the colony. The
  cheapest volume in the model is the only one that gives the pressurisation
  budget nothing back.
\end{enumerate}

Only one class of site in the model generates a rille, and therefore a tube.
This is a strong constraint and an intentional one; its consequences are
reported in Section~\ref{sec:results}.
